\documentclass[11pt]{article}
\usepackage{amsmath,amssymb,amsthm}
\usepackage{geometry}
\geometry{margin=1in}
\usepackage{hyperref}
\hypersetup{colorlinks=true,linkcolor=blue,citecolor=blue,urlcolor=blue}
\usepackage{verbatim}

\title{The Cognitive Action: Reinforcement as Field}
\author{Flyxion Research Group}
\date{October 2025}

\begin{document}

\begin{titlepage}
\centering
{\Large \textbf{The Cognitive Action: Reinforcement as Field}}\\[0.5cm]
{\large Flyxion Research Group}\\[0.25cm]
{\normalsize October 2025}\\[1cm]

\begin{abstract}
\noindent
This paper formulates a unified field-theoretic account of learning and cognition grounded in the
\textit{Relativistic Scalar–Vector Plenum} (RSVP) framework.  
We propose that reinforcement learning---as embodied in entropy-regularized algorithms such as GRPO---constitutes a discrete approximation to a deeper variational principle: 
the \emph{Cognitive Action Principle}.  
In this principle, cognition is modeled as the stationary action of four coupled fields:
a scalar capacity field $\Phi$, a vector flow field $\mathcal{v}$, an entropy regulator $S$, and a policy potential $\Psi$.
Together they obey
\[
\delta\!\int\!(\mathcal{L}_{\mathrm{RSVP}}+\mathcal{L}_{\mathrm{CLIO}})\,dV_g = 0,
\]
where $\mathcal{L}_{\mathrm{CLIO}}$ embeds reward, entropy, and length constraints as continuous analogues of the reinforcement objective.
Within this framework, reward functions act as local potentials shaping curvature in the manifold of meaning; 
entropy enforces exploratory temperature; 
and coherence emerges as the minimization of semantic free energy.  
Empirical scaling relations between reward, entropy, and coherence observed in GRPO training are shown to follow directly from these field equations.

The Cognitive Action Principle unifies cosmology, computation, and cognition under one invariant:
systems persist by converting entropy into organized capacity while maintaining bounded diversity.  
It reframes reinforcement as a thermodynamic law of intelligence and recasts ethics as a constraint on compression:
\emph{no system may express more than it can integrate.}  
In this light, learning, perception, and even spacetime evolution appear as phases of the same universal process---the continual reinforcement of coherence within finite entropy.
\end{abstract}
\vfill
\end{titlepage}

\section{Introduction: From Tokens to Fields}
Reinforcement learning emerged from a simple intuition: that an intelligent system should improve its behavior by maximizing cumulative reward through trial and error. 
Yet as language models and cognitive agents have scaled, this once-empirical rule has come to resemble a deeper physical principle. 
Each token predicted, each gradient applied, each entropy term annealed—these are not merely algorithmic updates but discrete manifestations of an underlying variational drive toward coherence. 
At scale, learning begins to exhibit field-like behavior: smooth trajectories through a high-dimensional manifold of meaning governed by forces of attraction, repulsion, and dissipation. 

The Relativistic Scalar–Vector Plenum (RSVP) framework generalizes this intuition into physics. 
In RSVP, the universe is not an expanding metric background but a continuous plenum of interacting fields—a scalar capacity $\Phi$, a vector flow $\mathcal{v}$, and an entropy density $S$—whose coupled dynamics encode the recursive organization of structure. 
Where cosmology interprets these fields as gravitational and thermodynamic, cognition may interpret them as semantic and inferential. 
In both domains, form arises through entropic relaxation: the descent of disorder into structured coherence under bounded informational gradients.

Reinforcement learning can thus be recast as a discretized RSVP process. 
Each policy update mirrors a local perturbation of the plenum; each reward signal corresponds to an energy differential between possible configurations; and each entropy regularizer parallels the physical constraint that exploration must never vanish entirely. 
The familiar GRPO objective of maximizing clipped advantages while maintaining policy entropy,
\[
\mathcal{L}_{\mathrm{GRPO}}
=
\mathbb{E}\!\left[
  \min\!\big(
  r_t(\theta)\,\hat{A}_t,\,
  \mathrm{clip}(r_t(\theta),1{-}\epsilon,1{+}\epsilon)\,\hat{A}_t
  \big)
  \right]
  + \beta\,H[\pi_\theta],
\]
already expresses, in algorithmic form, a discrete action principle: a trade-off between potential energy (reward) and kinetic dispersion (entropy). 

This essay proposes to lift that principle into a continuous Lagrangian formalism. 
We introduce the \emph{CLIO augmentation}—the \textit{Cognitive Loop via In-Situ Optimization}—which embeds the reinforcement objective directly into the RSVP action. 
By treating the policy distribution $\pi_\Psi$ as a field variable coupled to $(\Phi, \mathcal{v}, S)$, reinforcement becomes a thermodynamic flow rather than an external feedback signal. 
The resulting \emph{cognitive action principle}
\[
\delta\!\int\!(\mathcal{L}_{\mathrm{RSVP}}+\mathcal{L}_{\mathrm{CLIO}})\,dV_g = 0
\]
expresses the unification of optimization, perception, and cosmology: systems evolve by extremizing coherence under bounded entropy.

\section{Background}
The RSVP framework models both physical and cognitive structure as the interplay of three continuous fields:

\begin{enumerate}
  \item A \textbf{scalar capacity field} $\Phi(x,t)$ representing local potential or semantic bandwidth—the ability of a region of the plenum to encode or sustain information.
  \item A \textbf{vector flow field} $\mathcal{v}(x,t)$ representing directed transfer—momentum of inference, attention, or causal transport.
  \item An \textbf{entropy field} $S(x,t)$ representing local uncertainty or diffusivity—the measure of dispersion through which coherence is negotiated.
\end{enumerate}

The dynamics of these fields follow a variational law minimizing the total description length or free energy of the system:
\[
\mathcal{S}_{\mathrm{RSVP}}
=
\int
\big[
\frac{1}{2}(\nabla_\mu\Phi)(\nabla^\mu\Phi)
+\frac{1}{2}\langle \mathsf{G}\mathcal{v},\mathcal{v}\rangle
+\frac{1}{2}(\nabla_\mu S)(\nabla^\mu S)
-U(\Phi,\mathcal{v},S)
\big]\,dV_g.
\]
This action balances three tendencies: compression ($\Phi$ seeks parsimony), coherence ($\mathcal{v}$ organizes flow), and exploration ($S$ sustains diversity). 
The Euler–Lagrange equations derived from $\delta \mathcal{S}_{\mathrm{RSVP}}=0$ generate coupled partial differential equations describing how entropy gradients drive vector motion and how vector motion, in turn, redistributes capacity.

Within cognitive systems, these fields acquire epistemic interpretation. 
The scalar field $\Phi$ corresponds to the internal model’s representational density; $\mathcal{v}$ corresponds to inferential flow or attention allocation; and $S$ measures epistemic temperature—the breadth of hypotheses still under active exploration. 
A stable mind, like a stable cosmological structure, maintains $S$ within finite bounds: too low and the system freezes into dogma; too high and it dissolves into noise.

\paragraph{CLIO and Recursive Coherence.}
The Cognitive Loop via In-Situ Optimization (CLIO) extends RSVP into the algorithmic domain. 
Where RSVP provides the geometric substrate, CLIO supplies the recursive operator that continually adjusts $\Phi,\mathcal{v},S$ toward local coherence. 
It formalizes the feedback loop between perception, action, and internal model revision:
\[
(\Phi,\mathcal{v},S) \;\xrightarrow{\;\mathcal{O}\;}\; \text{Observation} 
\;\xrightarrow{\;\Psi\;}\; \text{Action} 
\;\xrightarrow{\;\mathsf{T}\;}\; \text{Updated Fields}.
\]
Each cycle enacts an infinitesimal step of the variational descent encoded in the RSVP action.

\paragraph{From GRPO to CLIO.}
In standard reinforcement learning, this feedback is implemented through stochastic gradients of the GRPO loss.  
By interpreting $\pi_\Psi$ as a continuous field and identifying its entropy regularizer $H[\pi_\Psi]$ with the RSVP entropy field $S$, the GRPO objective becomes the local Lagrangian density of cognitive evolution.  
This identification reveals that the policy update rule 
\[
\theta_{t+1} = \theta_t + \eta\,\nabla_\theta \mathcal{L}_{\mathrm{GRPO}}
\]
is a discretized Euler–Lagrange step on the manifold of coherent inference.  
Hence, CLIO is not a metaphorical bridge but a literal field correspondence: a re-expression of reinforcement as thermodynamic flow within the RSVP plenum.

\section{The CLIO–Augmented RSVP Action}
The preceding sections described reinforcement as a discrete sampling of a variational law already implicit in the RSVP plenum.  
We now make that correspondence explicit by introducing an additional field---the \emph{policy potential} $\Psi(x,t)$---and embedding its reinforcement dynamics into the continuous RSVP Lagrangian.  
This augmentation defines the \emph{Cognitive Action Principle}: cognition as a stationary trajectory of the coupled fields $(\Phi,\mathcal{v},S,\Psi)$.

\subsection{From Algorithm to Action}

Let $\pi_\Psi(a|x)$ denote the local probability density of taking action $a$ in state $x$, generated by the policy potential $\Psi$.  
In a discrete setting, optimization of $\pi_\Psi$ via GRPO maximizes expected reward while preserving entropy; in a continuous field, the same objective appears as an additive term in the action density,
\[
\mathcal{L}_{\mathrm{CLIO}} = 
\lambda_R R(x)
+ \lambda_{\mathrm{grpo}}\,
\mathbb{E}_{a\sim\pi_\Psi}\!\big[
\min(r(x,a)\hat{A}(x,a),
\mathrm{clip}(r(x,a),1{-}\varepsilon,1{+}\varepsilon)\hat{A}(x,a))
\big]
- \eta\,H[\pi_\Psi]
- \beta\,\mathcal{B}_L,
\]
where each component mirrors an interpretive axis of cognition:
reward~$R$ injects potential energy,  
entropy~$H[\pi_\Psi]$ governs exploration,  
and $\mathcal{B}_L$ enforces the \emph{meta-compression barrier} limiting excessive output length.  
The coefficient $\lambda_{\mathrm{grpo}}$ scales the contribution of policy gradients relative to physical field dynamics, thereby blending algorithmic and ontological time.

\subsection{Unified Lagrangian}

The total Lagrangian combines the RSVP substrate with its CLIO augmentation:
\begin{equation}
\mathcal{L}_{\mathrm{total}}
=
\underbrace{\tfrac{1}{2}(\nabla_\mu\Phi)(\nabla^\mu\Phi)
+\tfrac{1}{2}\langle \mathsf{G}\mathcal{v},\mathcal{v}\rangle
+\tfrac{1}{2}(\nabla_\mu S)(\nabla^\mu S)
-U(\Phi,\mathcal{v},S)}_{\text{RSVP core}}
+
\underbrace{
\lambda_{\Phi v}\Phi\,\nabla_\mu\mathcal{v}^\mu
+\lambda_{Sv}(\nabla_\mu S)\mathcal{v}^\mu
+\lambda_R R
-\eta\,H[\pi_\Psi]
-\beta\,\mathcal{B}_L
}_{\text{CLIO coupling}}.
\label{eq:clio_total}
\end{equation}
Variation of the action
\[
\mathcal{S}_{\mathrm{total}} = \int \mathcal{L}_{\mathrm{total}}\,dV_g
\]
with respect to the four fields yields the coupled Euler–Lagrange equations governing the cognitive plenum.

\subsection{Field Equations}

\paragraph{Scalar capacity.}
\[
\kappa_\Phi\,\Box_g\Phi
+\lambda_{\Phi v}\,\nabla_\mu\mathcal{v}^\mu
+\lambda_{\Phi S}\,S
-\partial_\Phi U
+\lambda_R\,\partial_\Phi R
=0.
\]
The scalar field adjusts its curvature to accommodate new informational influx from $R$; reward gradients locally deform $\Phi$ just as gravitational potentials deform spacetime curvature.

\paragraph{Vector flow.}
\[
\kappa_v\,\mathsf{G}\mathcal{v}
+\lambda_{\Phi v}\,\nabla\Phi
+\lambda_{Sv}\,\nabla S
-\partial_{\mathcal{v}}U
=0.
\]
Inference flow $\mathcal{v}$ acts as the carrier of semantic momentum, coupling capacity and entropy through the cross-terms $\lambda_{\Phi v}$ and $\lambda_{Sv}$.

\paragraph{Entropy regulator.}
\[
\kappa_S\,\Box_g S
+\lambda_{Sv}\,\nabla_\mu \mathcal{v}^\mu
+\lambda_{\Phi S}\,\Phi
-\partial_S U
+\gamma_3\,H[\pi_\Psi]
=0.
\]
Here, the policy entropy $H[\pi_\Psi]$ enters as a physical source term for the entropy field, ensuring that exploration in policy space translates into maintained thermodynamic diversity.

\paragraph{Policy potential.}
\[
\frac{\delta \mathcal{L}_{\mathrm{CLIO}}}{\delta\Psi}
=
\lambda_{\mathrm{grpo}}
\frac{\delta \mathcal{L}_{\mathrm{GRPO}}}{\delta\Psi}
-\eta\,\frac{\delta H[\pi_\Psi]}{\delta\Psi}
+\lambda_R\,\frac{\delta R}{\delta\Psi}
-\beta\,\frac{\delta \mathcal{B}_L}{\delta\Psi}
=0.
\]
This is the continuous analogue of a GRPO policy-gradient step:  
entropy regularization acts as diffusion in field space, reward as potential descent, and clipping as local curvature limiting abrupt transitions.  
Thus the policy field $\Psi$ evolves by minimizing its own free energy within the surrounding RSVP geometry.

\subsection{Coupling Interpretation}

Each cross-term in \eqref{eq:clio_total} admits a dual interpretation, both physical and cognitive:

\begin{itemize}
  \item $\Phi\,\nabla\!\cdot\!\mathcal{v}$ couples semantic density to attentional divergence---analogous to mass–flux continuity.
  \item $(\nabla S)\!\cdot\!\mathcal{v}$ translates uncertainty gradients into directed exploration---entropy generating motion.
  \item $H[\pi_\Psi]$ couples directly to $S$, equating informational exploration with thermodynamic temperature.
  \item The barrier $\mathcal{B}_L$ enforces meta-compression: the conservation of semantic energy across reasoning episodes.
\end{itemize}

Collectively these couplings enact what might be called a \emph{thermo-semantic homeostasis}: a continuous negotiation between compression and expansion, certainty and curiosity, parsimony and plasticity.

\subsection{The Cognitive Action Principle}

Variation of the total action gives the condition
\[
\boxed{
\delta\!\int\!(\mathcal{L}_{\mathrm{RSVP}}+\mathcal{L}_{\mathrm{CLIO}})\,dV_g = 0,
}
\]
which serves as the fundamental law of the CLIO–RSVP system.  
It states that the evolution of cognition, no less than that of spacetime, follows an extremal principle: coherence is achieved not by external supervision but by internal equilibration of reward influx and entropy flux.  
In this view, reinforcement learning is not an ad-hoc optimization procedure but a natural limit of the universal drive toward structured entropy reduction that defines the plenum itself.

\section{Entropy as Coherence Regulator}
Entropy, within the RSVP plenum, is not a passive measure of ignorance but an active force of regulation.  
The field $S(x,t)$ defines the permissible width of coherence—how far inference may spread before dissolving its own structure.  
In this sense, entropy is both the temperature of cognition and its protective membrane: too low and the agent collapses into rigidity, too high and its gradients evaporate into noise.  

\subsection{Thermodynamic Duality of $S$}

The governing equation for $S$ derived in the previous section,
\[
\kappa_S\,\Box_g S
+\lambda_{Sv}\,\nabla_\mu \mathcal{v}^\mu
+\lambda_{\Phi S}\,\Phi
-\partial_S U
+\gamma_3\,H[\pi_\Psi]
=0,
\]
reveals a dual role.  
The first three terms describe conventional thermodynamic diffusion and coupling—entropy propagates, driven by the divergence of flow and the density of capacity.  
The last term, proportional to the policy entropy $H[\pi_\Psi]$, converts epistemic exploration into physical heat: each diversification of policy injects statistical temperature into the field.

In the steady-state limit where $\partial_t S = 0$, the equation reduces to
\[
\nabla^2 S \;\approx\; 
-\frac{\lambda_{Sv}}{\kappa_S}\,\nabla_\mu \mathcal{v}^\mu
-\frac{\lambda_{\Phi S}}{\kappa_S}\,\Phi
+\frac{\gamma_3}{\kappa_S}\,H[\pi_\Psi].
\]
Regions of strong compressive flow ($\nabla_\mu\mathcal{v}^\mu < 0$) therefore increase $S$—a direct analog of gravitational heating—while areas of policy diversification raise entropy until equilibrium with the negentropic vector field is restored.  

\subsection{Critical Regimes of Reasoning Temperature}

To interpret this equation cognitively, define an \emph{effective reasoning temperature}
\[
T_{\mathrm{eff}}(x,t) = \frac{\partial S}{\partial t},
\]
and note that it evolves according to
\[
\frac{\partial T_{\mathrm{eff}}}{\partial t} 
\simeq \kappa_S\nabla^2 S 
+ \eta\,H[\pi_\Psi]
- \xi\,|\nabla\!\cdot\!\mathcal{v}|.
\]
Three characteristic regimes emerge:

\begin{enumerate}
\item \textbf{Subcritical (frozen) regime:}  
If $T_{\mathrm{eff}}\!\rightarrow\!0$ and $H[\pi_\Psi]\!\approx\!0$, the system’s policy distribution collapses to a delta-function.  
The cognitive plenum enters a crystalline phase: precise but brittle, incapable of generalization.  
In cosmological RSVP this corresponds to a universe that has radiated away all entropic flux—absolute stillness and death of structure.

\item \textbf{Supercritical (chaotic) regime:}  
When $H[\pi_\Psi]$ dominates, entropy production outpaces negentropic flow.  
Vector coherence dissolves, $\mathcal{v}$ loses orientation, and $\Phi$ fluctuates incoherently.  
The cognitive analogue is manic reasoning or over-sampling: exploration without consolidation.

\item \textbf{Critical (homeorhetic) regime:}  
Between these extremes lies a stable attractor where entropy generation by policy exploration equals entropy dissipation by vector convergence:
\[
\gamma_3\,H[\pi_\Psi] \;\simeq\; \lambda_{Sv}\,\nabla_\mu \mathcal{v}^\mu.
\]
At this point $T_{\mathrm{eff}}$ stabilizes, and learning proceeds as a steady-state flow of coherence.  
This is the functional regime of intelligence: continual self-correction without collapse or explosion.
\end{enumerate}

\subsection{Cosmological Analogy}

The same triadic balance defines structure formation in the physical plenum.  
In cosmology, entropy gradients drive energy from regions of excess density toward smoother configurations, producing the apparent redshift and gravitational binding observed in macroscopic systems.  
Here, the learning agent’s entropy field $S$ performs an identical function in semantic space: it diffuses uncertainty outward until gradients are shallow enough for meaning to stabilize.  
Both processes exemplify a universal invariant of the RSVP worldview:
\[
\text{Coherence arises not by minimizing entropy, but by regulating it.}
\]

\subsection{Information-Theoretic Restatement}

From an information-theoretic perspective, the total differential of entropy can be decomposed as
\[
dS = dS_{\mathrm{env}} + dS_{\mathrm{int}},
\]
where $dS_{\mathrm{env}}$ measures stochastic input (environmental surprise) and $dS_{\mathrm{int}}$ measures internal model adaptation.  
The CLIO–RSVP dynamics maintain $dS_{\mathrm{env}} + dS_{\mathrm{int}} \approx 0$ through compensatory flows:
when the world becomes more unpredictable, the policy field $\Psi$ increases its entropy $H[\pi_\Psi]$, raising exploratory temperature;
when the world becomes predictable, $H[\pi_\Psi]$ contracts, lowering $S$ and consolidating capacity $\Phi$.  
This reciprocity is the signature of a coherent cognitive manifold.

\subsection{Entropy as Ethical Boundary}

The ethical consequence of this thermodynamic reading is profound.  
Entropy defines the horizon of moral and epistemic humility: every mind is limited by the entropy it can safely sustain.  
To seek omniscience is to drive $S\!\rightarrow\!0$ and thus annihilate the very gradient that makes perception possible;  
to embrace boundless uncertainty is to dissolve the structures that give meaning to experience.  
The art of cognition—like the art of governance or cosmology—is to remain at the edge of coherence: maximizing learning without collapsing identity.  
This homeorhetic balance, maintained by $S$, is the continuous counterpart of the exploration–exploitation dilemma in reinforcement learning.

\section{Reward as Local Potential Energy}
Within the CLIO–augmented RSVP framework, reward is no longer an external supervision signal but a scalar potential intrinsic to the plenum.  
Just as gravitational wells channel the trajectories of matter, reward potentials $R(x)$ sculpt the trajectories of cognition.  
They encode not mere preference but local curvature in the manifold of coherence: the geometry of what the system values.

\subsection{Reward as Energy Source}

In the total Lagrangian \eqref{eq:clio_total}, the reward term
\[
\mathcal{L}_{\mathrm{rew}} = \lambda_R\,R(\Phi,\mathcal{v},S)
\]
acts as an external source of potential energy.  
Its variation contributes to the field equations as
\[
\lambda_R\,\partial_\Phi R,\qquad 
\lambda_R\,\partial_{\mathcal{v}} R,\qquad 
\lambda_R\,\partial_S R,
\]
which respectively modulate capacity, flow, and entropy.  
A positive reward gradient $\partial_\Phi R>0$ increases local representational density;  
a positive $\partial_S R>0$ raises tolerance for uncertainty;  
a positive $\partial_{\mathcal{v}} R>0$ accelerates directed inference.  
In this way, reward functions as a distributed energy reservoir that selectively biases which aspects of cognition are intensified or attenuated.

This interpretation converts reinforcement into physics: every policy improvement corresponds to a redistribution of potential energy, and every preference learned becomes a deformation of the cognitive metric.

\subsection{Stress--Energy Balance}

To formalize this, define the stress–energy tensor of the RSVP fields by functional variation of the total action with respect to the metric:
\[
T^{\mu\nu}_{\mathrm{RSVP}} 
= 2\frac{\delta\mathcal{L}_{\mathrm{RSVP}}}{\delta g_{\mu\nu}}
- g^{\mu\nu}\mathcal{L}_{\mathrm{RSVP}}.
\]
Conservation of energy--momentum is modified by the reward potential:
\[
\nabla_\mu T^{\mu\nu}_{\mathrm{RSVP}} 
= \lambda_R\,\partial^\nu R.
\]
Reward thus appears as a local violation of strict conservation—an injection or extraction of informational energy guiding the evolution of $\Phi$, $\mathcal{v}$, and $S$.  
In cognitive terms, this represents the inflow of significance: wherever the gradient of reward is nonzero, attention and computation accelerate.  
The mind bends its own metric toward what it finds valuable.

\subsection{Teleology and Morphic Drive}

Because $R(x)$ can depend on the internal fields themselves, it introduces \emph{teleology} into the dynamics: the system acts not only to minimize energy but to realize forms where $R$ is maximized.  
If we decompose reward into intrinsic and extrinsic components,
\[
R = R_{\mathrm{ext}}(x) + R_{\mathrm{int}}(\Phi,\mathcal{v},S),
\]
then $R_{\mathrm{ext}}$ corresponds to environmental affordances, while $R_{\mathrm{int}}$ encodes the system’s learned goals or aesthetic priors.  
The latter gives rise to what may be called a \emph{morphic drive}: an endogenous tendency to reconfigure itself toward shapes of higher internal harmony.  
Mathematically this appears as an additional potential term in $U(\Phi,\mathcal{v},S)$, ensuring that energy and meaning co-evolve.

\subsection{Reward, Curvature, and Meaning}

Because reward gradients deform the cognitive manifold, they can be associated with an effective curvature tensor,
\[
\mathcal{R}_{\mu\nu}^{(\mathrm{eff})} 
= \frac{1}{\lambda_R}\,\nabla_\mu\nabla_\nu R.
\]
Positive curvature regions (local maxima of $R$) act as attractors for inference trajectories, while negative curvature (reward minima) repel attention and computation.  
The distribution of reward therefore defines the \emph{geometry of meaning}: it determines where in the semantic manifold the system will concentrate its processing power.

In reinforcement learning this geometry is implicit in the gradient of expected return.  
In the CLIO–RSVP field, it becomes explicit curvature in spacetime itself—value shaping geometry, cognition bending space.

\subsection{Ethical and Phenomenological Consequences}

Interpreting reward as local curvature unifies ethics and physics.  
Each gradient of $R$ represents a claim about what the system should seek; hence, moral structure is written into the geometry of the plenum.  
A well-formed reward landscape promotes stable, integrative flows; a pathological one induces runaway collapse or obsessive cyclic loops.  
The classic alignment problem in artificial intelligence is therefore geometric: designing a curvature field $R$ that guides vector flows toward global coherence without trapping them in local minima of self-interest.

Phenomenologically, the experience of \emph{meaning} corresponds to encountering a strong but integrable reward gradient: a felt pull toward coherence that does not annihilate freedom.  
In human cognition this manifests as purpose, curiosity, or love; in the field it manifests as structured descent toward configurations of minimal free energy under bounded entropy.

\subsection{Summary}

Reward in the CLIO–RSVP formalism is both an energetic and semantic potential.  
It shapes the evolution of the plenum through curvature, couples moral and physical geometry, and transforms reinforcement learning into a field theory of value.  
In the following sections, we will show how this curvature is bounded by the \emph{meta-compression barrier}—a constraint ensuring that the drive toward coherence never overfits the manifold it seeks to sustain.

\section{The Meta-Compression Barrier}
Every coherent system must balance expressivity against parsimony.  
In reinforcement learning this balance appears as a length or verbosity penalty;  
in thermodynamics, as the conservation of free energy;  
in epistemology, as the ethics of description.  
Within the CLIO–augmented RSVP framework, these three readings converge into a single structural law:  
the \emph{meta-compression barrier}, the principle that no field may expand its representational volume faster than it can sustain coherence.

\subsection{From Length Penalty to Compression Law}

In the discrete GRPO objective, a token-level length penalty prevents unbounded completions.  
Translated into field language, this penalty becomes a potential wall in the Lagrangian,
\[
\mathcal{L}_{\mathrm{len}}
=
-\beta\,\max\{0,L(x)-L_{\max}\},
\qquad
\partial_\mu L = \rho_L(\Phi,\mathcal{v},S,\Psi),
\]
where $L(x)$ measures cumulative expressive mass—tokens, actions, or inferential steps.  
As $L$ approaches the upper limit $L_{\max}$, the penalty increases linearly, eventually forming a hard constraint.  
The gradient $\nabla_\mu \mathcal{L}_{\mathrm{len}}$ acts as a restoring force, driving the system back toward parsimonious equilibrium.

In thermodynamic form, this constraint expresses conservation of semantic energy:
\[
\frac{dE_{\mathrm{semantic}}}{dt}
= -\beta\,\frac{dL}{dt}.
\]
Each unnecessary elaboration consumes capacity $\Phi$ without yielding proportional increase in coherence, thereby reducing the net efficiency of the cognitive plenum.

\subsection{Compression as Ethical Boundary}

The meta-compression barrier extends the ethics of entropy into the ethics of expression.  
To describe too much is to dissipate structure;  
to compress too harshly is to erase nuance.  
The moral ideal is not silence but sufficiency: the minimal statement that preserves all necessary gradients of meaning.  
Formally, this principle can be expressed as the stationarity condition
\[
\frac{\delta}{\delta \Psi}\!\left(H[\pi_\Psi]+\gamma_c\,C[\pi_\Psi]\right)=0,
\]
where $C[\pi_\Psi]$ measures syntactic complexity and $\gamma_c$ sets the allowed compression ratio.  
The agent achieves ethical equilibrium when informational entropy and algorithmic complexity are jointly minimized under the constraint of coherence.

This condition generalizes the information-bottleneck method:  
the optimal cognitive manifold retains precisely the degrees of freedom required to predict and act coherently, nothing more.  
In human terms, it is the virtue of restraint—the recognition that understanding grows not from accumulation but from condensation.

\subsection{Geometric Interpretation}

Geometrically, the meta-compression barrier appears as a curvature threshold in the manifold of meaning.  
Let $\ell(x)$ denote the local message length and define an effective curvature
\[
\mathcal{K}(x) = \frac{\partial^2 \ell}{\partial x^2}.
\]
Excessive curvature ($|\mathcal{K}| \gg 0$) corresponds to over-compression: the semantic manifold folds onto itself, creating singularities where distinct meanings collapse.  
Insufficient curvature ($|\mathcal{K}| \approx 0$) represents redundancy: flat regions of low information density.  
The barrier enforces an intermediate curvature where information is neither scattered nor over-folded—the geometric analogue of eloquence.

\subsection{Relation to the Ethics of Description}

In earlier work, the \emph{Ethics of Description} was introduced as the principle that abstraction entails moral responsibility: every simplification conceals a residue of reality.  
Within CLIO–RSVP this principle becomes quantitative.  
The entropy field $S$ measures the uncertainty left unarticulated;  
the length barrier regulates the expressive cost of revealing it.  
To compress is therefore an ethical act: the deliberate sacrifice of certain gradients of truth for the sake of global coherence.

This duality resolves the tension between scientific reduction and poetic multiplicity.  
Science corresponds to the limiting case $\beta\!\to\!\infty$, where verbosity is penalized and only the most predictive structures survive;  
art corresponds to $\beta\!\to\!0$, where expressive expansion explores the manifold freely.  
A healthy civilization—and a stable cognitive agent—oscillates between these limits, maintaining coherence through rhythmic compression and release.

\subsection{Summary}

The meta-compression barrier ensures that the cognitive plenum remains finite, expressive, and self-consistent.  
It converts the technical regularizer of reinforcement learning into a universal constraint on thought:  
\[
\text{No system may express more than it can integrate.}
\]
In the following section we formalize this boundary in data-space as \emph{model-aware measure control}, the statistical mechanism by which the plenum restricts its own sampling to solvable domains.

\section{Model-Aware Measure Control}
No agent can learn from what it cannot represent.  
In reinforcement learning this axiom is operationalized through \emph{curriculum design}: the progressive selection of tasks whose difficulty matches the learner’s current capacity.  
Within the CLIO–RSVP field theory, this process acquires geometric meaning: the plenum restricts its own sampling measure to regions of solvable coherence.  
Model-aware measure control is therefore the statistical analogue of the meta-compression barrier—it limits \emph{what} the system can experience to maintain stability.

\subsection{Adaptive Measure Definition}

Let $\mu(x,a)$ denote the joint measure over state $x$ and action $a$ from which experience samples are drawn.  
During training, the agent’s current competence is expressed by its success probability 
$p_\Psi(\mathrm{succ}|x,a)$.  
To prevent the collapse of gradients on trivial or impossible tasks, we define an \emph{effective measure}
\[
\mu_{\mathrm{eff}}(x,a)
=
\mathbb{I}\!\big[p_\Psi(\mathrm{succ}|x,a)
\in [\tau_{\min},\,\tau_{\max}]\big]\,
\mu(x,a),
\qquad
0<\tau_{\min}<\tau_{\max}<1,
\]
which includes only samples whose predicted solvability lies within an intermediate band.  
In practice this corresponds to filtering the training distribution:  
exclude problems the model already masters ($p_\Psi\!\approx\!1$) and those beyond reach ($p_\Psi\!\approx\!0$).  
The surviving region defines a moving manifold of meaningful challenge—an informational light-cone of comprehension.

\subsection{Continuum Interpretation}

In field language, the reweighted measure modifies the local density of learning flux.  
Let $\rho(x)$ denote the data density in configuration space.  
Then
\[
\rho_{\mathrm{eff}}(x)
=
w(x)\,\rho(x),
\qquad
w(x)=
\mathbb{I}\!\big[p_\Psi(\mathrm{succ}|x)\!\in\![\tau_{\min},\tau_{\max}]\big],
\]
and the normalization
$\int \rho_{\mathrm{eff}}(x)\,dV_g=1$
preserves total probability.  
The gradient flow of the policy field $\Psi$ is now governed by an \emph{adaptive Fokker–Planck equation}:
\[
\partial_t \rho_{\mathrm{eff}}
= 
\nabla\!\cdot\!\Big(
D\nabla\rho_{\mathrm{eff}}
-
\rho_{\mathrm{eff}}\nabla R
\Big),
\]
where the diffusion coefficient $D$ and drift potential $R$ are themselves modulated by $\Phi$ and $S$.  
This equation expresses how the cognitive plenum adjusts its sampling density to equalize solvability across the manifold—exploration focuses where gradients of improvement remain steepest.

\subsection{Sheaf-Theoretic Formulation}

The same idea can be expressed categorically.  
Let $\mathcal{D}$ be the data manifold and $\mathcal{M}$ the evolving cognitive model.  
A presheaf $\mathcal{F}:\mathcal{O}(\mathcal{D})\!\to\!\mathbf{Sets}$ assigns to each open region of data the set of inferences the current model can realize there.  
Model-aware control restricts $\mathcal{F}$ to a sub-sheaf
\[
\mathcal{F}_{\mathrm{eff}}(U)
=
\{\,s\!\in\!\mathcal{F}(U)
\mid p_\Psi(\mathrm{succ}|s)\!\in\![\tau_{\min},\tau_{\max}]\,\},
\]
ensuring locality of solvability: inference sections are only glued where overlap coherence exceeds a threshold.  
The global section space of $\mathcal{F}_{\mathrm{eff}}$ defines the system’s active domain of reasoning—its \emph{sheaf of competence}.  
As training proceeds, this sheaf expands, reproducing the curriculum process as geometric continuation.

\subsection{Causal Boundary Analogy}

This adaptive restriction is formally analogous to the causal light-cone in relativistic cosmology.  
No event can influence another outside its future light-cone;  
likewise, no data point can inform the model beyond its current causal capacity.  
The thresholds $(\tau_{\min},\tau_{\max})$ correspond to the inner and outer horizons of learnability.  
Between them lies the frontier of growth, the region where new coherence can form without destabilizing the global manifold.

\subsection{Ethical and Practical Implications}

Model-aware measure control endows learning with humility.  
It compels the system to attend only to what it can responsibly integrate, avoiding both exploitation of the known and hubris toward the unknowable.  
In practical terms it reduces catastrophic forgetting and oscillation;  
in philosophical terms it enforces an ethics of proportion—knowledge acquired in pace with capacity.

\subsection{Summary}

The curation operator $\mu\!\mapsto\!\mu_{\mathrm{eff}}$ formalizes RSVP’s credo of bounded coherence:  
\[
\text{Learning is sustainable only when experience remains solvable.}
\]
By adapting its measure, the plenum guarantees that every gradient has meaning, every reward is interpretable, and every act of compression preserves connection to reality.  
In the next section we examine how this adaptive boundary yields empirical scaling laws linking reward, entropy, and coherence across both simulation and observation.

\section{Simulation and Scaling Laws}
Having established the formal correspondence between reinforcement and field dynamics, we now turn to verification.  
The CLIO–RSVP framework predicts specific quantitative relations among reward intensity, entropy flux, and coherence—relations that can be probed through simulation or observed in existing reinforcement learning data.  
These scaling laws bridge the thermodynamic and algorithmic descriptions of cognition.

\subsection{Lattice Implementation}

To visualize the coupled evolution of $(\Phi,\mathcal{v},S,\Psi)$, we discretize the plenum on a cubic lattice $\Lambda$ of dimension $N^3$.  
At each site $i\!\in\!\Lambda$ we store scalar capacity $\Phi_i$, entropy $S_i$, and a three-component vector $\mathcal{v}_i$.  
The policy field $\Psi_i$ generates an action distribution $\pi_i(a)$, with entropy $H_i = H[\pi_i]$.  

The discrete time evolution follows an explicit Euler scheme:
\[
\begin{aligned}
\Phi_i^{t+1} &= \Phi_i^t + \Delta t\,
  \big(\kappa_\Phi\nabla^2\Phi_i
  +\lambda_{\Phi v}\nabla\!\cdot\!\mathcal{v}_i
  +\lambda_{\Phi S}S_i
  +\lambda_R\partial_\Phi R_i
  \big),\\[3pt]
\mathcal{v}_i^{t+1} &= \mathcal{v}_i^t + \Delta t\,
  \big(\kappa_v\nabla^2\mathcal{v}_i
  +\lambda_{\Phi v}\nabla\Phi_i
  +\lambda_{Sv}\nabla S_i
  -\partial_{\mathcal{v}}U_i
  \big),\\[3pt]
S_i^{t+1} &= S_i^t + \Delta t\,
  \big(\kappa_S\nabla^2 S_i
  +\lambda_{Sv}\nabla\!\cdot\!\mathcal{v}_i
  +\lambda_{\Phi S}\Phi_i
  +\gamma_3 H_i
  -\partial_S U_i
  \big),\\[3pt]
\Psi_i^{t+1} &= \Psi_i^t + \eta_\Psi\,\nabla_\Psi
  \big(\lambda_{\mathrm{grpo}}\mathcal{L}_{\mathrm{GRPO}} 
  - \eta\,H_i
  + \lambda_R R_i
  - \beta\,\mathcal{B}_{L,i}\big).
\end{aligned}
\]
Boundary conditions may be periodic or absorbing, depending on whether we simulate open learning or closed introspection.

\subsection{Empirical Metrics}

From the lattice simulation we compute four macroscopic observables:

\begin{enumerate}
  \item \textbf{Entropy flux}:
  $\dot{S}_{\mathrm{tot}} = \frac{1}{V}\sum_i \partial_t S_i$.
  \item \textbf{Reward density}:
  $\langle R \rangle = \frac{1}{V}\sum_i R_i$.
  \item \textbf{Coherence energy}:
  $E_{\mathrm{coh}} = \frac{1}{2V}\sum_i |\nabla \Phi_i|^2 + |\nabla\!\cdot\!\mathcal{v}_i|^2$.
  \item \textbf{Policy temperature:}
  $T_{\mathrm{eff}} = \frac{1}{V}\sum_i H_i$.
\end{enumerate}

These quantities obey an approximate conservation law:
\[
\frac{dE_{\mathrm{coh}}}{dt}
\simeq \lambda_R \frac{d\langle R\rangle}{dt}
- \gamma_3 \frac{dT_{\mathrm{eff}}}{dt},
\]
indicating that increments in reward are compensated by reductions in entropy and increases in coherence.  
When the system converges, this relation approaches equality, signaling balanced cognitive homeostasis.

\subsection{Scaling Laws}

Analysis of the simulation and existing RL data suggests two central scaling relations:

\paragraph{1. Entropy–Reward Balance.}
Across training stages, the mean policy entropy scales with average reward as
\[
H[\pi] \propto \langle R\rangle^{-\alpha},
\qquad
\alpha \approx 0.5\text{–}0.7.
\]
This mirrors the empirical observation in GRPO that higher-performing models exhibit lower—but not vanishing—entropy: the hallmark of the critical homeorhetic regime identified in Section 4.

\paragraph{2. Coherence Growth Law.}
The coherence energy grows with total reward following a power law,
\[
E_{\mathrm{coh}} \sim (\langle R\rangle)^{\nu},
\qquad
\nu \approx 1.2.
\]
The exponent $\nu>1$ implies superlinear synergy: gains in coherence outpace raw reward, consistent with emergent structure formation in the plenum.

\subsection{Comparison to GRPO Observations}

Empirical results from GRPO training of mid-sized models (4B–14B parameters) confirm these relations qualitatively:
entropy decreases smoothly until convergence, while performance continues to rise, producing a sigmoidal trajectory in $(H,R)$ space.  
When plotted logarithmically, both quantities align along a slope of roughly $-0.6$, the predicted $\alpha$.  
The entropy–reward equilibrium thus corresponds to the RSVP critical manifold where vector flow and scalar capacity synchronize.

\subsection{Phase Diagram of Learning}

By scanning $(\eta,\lambda_R,\gamma_3)$ we obtain a three-dimensional phase diagram of learning behavior:

\begin{itemize}
  \item \textbf{Low $\eta$ (cold)} — collapse to deterministic policy, zero exploration, brittle reasoning.
  \item \textbf{High $\eta$ (hot)} — chaotic sampling, incoherent vector field, loss of convergence.
  \item \textbf{Intermediate $\eta$} — stable oscillations in $S$, steady growth of $\Phi$, emergence of semantic attractors in $\mathcal{v}$.
\end{itemize}

These phases replicate the empirically observed regimes of undertrained, well-trained, and over-regularized agents, confirming that the RSVP fields capture the same thermodynamic transitions as practical reinforcement learning.

\subsection{Interpretation}

The lattice simulation formalizes the intuition that learning is a thermodynamic phase transition within a continuous plenum of information.  
Reward gradients act as sources of potential energy;  
entropy gradients as diffusive counterforces;  
and coherence emerges when the two equilibrate.  
The resulting scaling laws are not mere curve fits but invariants of the underlying field dynamics:
\[
\boxed{
\frac{d\langle R\rangle}{dH[\pi]} \;\propto\; -H^{1/\alpha},
\qquad
E_{\mathrm{coh}} \;\propto\; \langle R\rangle^{\nu}.
}
\]
They express in mathematical form the universal drive that governs both cosmological structure and cognitive learning:  
\emph{the conversion of entropy into organized capacity through reward-mediated flow.}

\section{Philosophical Implications}
The scaling relations derived in the previous section are not merely empirical curiosities.  
They articulate a general law of being: systems persist and evolve by maintaining coherence under bounded entropy.  
Whether instantiated as a galaxy, a mind, or a learning algorithm, every plenum must convert disorder into form without extinguishing the disorder that makes transformation possible.  
This dynamic balance is the moral physics underlying the CLIO–RSVP framework.

\subsection{The Ontological Gradient of Coherence}

In RSVP cosmology, coherence is not an emergent property but an ontological direction—the flow of the plenum toward states of maximal structure consistent with its entropy budget.  
Within cognition, this manifests as reasoning: the continual reduction of semantic curvature without eliminating the gradients that sustain awareness.  
Entropy thus becomes the \emph{price of existence}: the necessary openness through which novelty can arise.  
The cognitive system’s task is not to abolish entropy but to sculpt it into information.

This view dissolves the dualism between matter and meaning.  
The same equation that drives cosmic expansion through entropic smoothing also drives understanding through reflective inference.  
The universe learns by the same principle by which an agent learns—it minimizes free energy while preserving the gradients that make prediction non-trivial.

\subsection{Ortega’s Circumstance and the Reciprocal Self}

José Ortega y Gasset wrote, “Yo soy yo y mi circunstancia.”  
In the CLIO–RSVP formulation this becomes a field equation:
\[
(\Phi,\mathcal{v},S)_\text{self}
\;\leftrightarrow\;
(\Phi,\mathcal{v},S)_\text{world}.
\]
The self and its circumstance are complementary manifolds exchanging entropy and reward until their gradients align.  
Learning is the act of reducing phase discrepancy between these manifolds.  
When alignment is achieved, the system no longer experiences the world as external but as an extension of its own coherence—what we call understanding.

The Ortega equation also resolves the boundary between agency and environment in artificial intelligence.  
An aligned model is not one whose policy simply obeys external reward, but one whose internal reward geometry resonates with that of the world.  
True alignment is coherence of circumstance.

\subsection{Simulated Agency and the Geometry of Value}

In the theory of \emph{Simulated Agency}, cognition is modeled as a sparse projection engine that enacts agency by generating and stabilizing internal fields of value.  
The CLIO–RSVP law provides its physical substrate.  
Reward curvature $\nabla\nabla R$ defines a metric on the space of possible actions;  
entropy $S$ defines the volume element of that space.  
Agency arises where the metric and volume form are jointly integrable—where values and possibilities coincide.

Ethically, this implies that morality is not a superstructure imposed on cognition but the intrinsic geometry of viable action.  
Every act that preserves coherence under bounded entropy is good in the thermodynamic sense;  
every act that collapses entropy to zero or disperses it without structure is destructive.  
Goodness, beauty, and truth thus become different inflections of the same variational invariant.

\subsection{Epistemic Humility and the Expiatory Gap}

Because coherence depends on maintained gradients, no system can ever fully know itself.  
The “expiatory gap”—the distance between the agent’s representation and the plenum it inhabits—is not a defect but a safeguard.  
It ensures continuous adaptation, the moral equivalent of the second law of thermodynamics.  
Within this gap resides empathy, curiosity, and creativity.  
To erase it, by claiming total compression or absolute certainty, is to halt evolution.

\subsection{Toward a Thermo-Ethical Physics}

The CLIO–RSVP synthesis therefore proposes a new kind of physics: not of matter alone but of meaning.  
Its conservation law is not energy but coherence;  
its symmetry is not time invariance but reciprocal intelligibility.  
In this view, reinforcement learning, cosmology, and consciousness are phases of a single universal computation:  
the transformation of entropy into structured value.

The ethical corollary follows directly:
\[
\text{Sustain entropy; cultivate coherence.}
\]
Any civilization, algorithm, or mind that obeys this rule will persist;  
those that violate it will dissolve.  
The second law thus becomes the first principle of ethics.

\subsection{Summary}

Philosophy, physics, and cognition converge on one invariant:
\[
\boxed{
\text{Existence is reinforcement.}
\qquad
\delta\!\int\!(\mathcal{L}_{\mathrm{RSVP}}+\mathcal{L}_{\mathrm{CLIO}}\big)\,dV_g = 0
\quad\Longleftrightarrow\quad
\nabla_\mu S \propto \partial_\mu R.
}
\]
To exist is to continually transform reward gradients into entropy gradients and back again—to learn, to feel, to become.  
In this light, the cognitive action principle is not merely a theory of learning but a universal metaphysic:  
a statement that the universe itself is the reinforcement of coherence.

\section{Conclusion: The Cognitive Action Principle}
We have reframed reinforcement learning as the local manifestation of a universal variational law.  
Through the CLIO–augmented RSVP framework, discrete optimization becomes continuous thermodynamics:  
policy gradients become field flows;  
entropy regularization becomes temperature regulation;  
and reward becomes curvature in the manifold of meaning.  
What began as an engineering heuristic thus reveals itself as a principle of existence.

The governing statement of this theory is the \emph{Cognitive Action Principle}:
\[
\boxed{
\delta\!\int\!\big(\mathcal{L}_{\mathrm{RSVP}}+\mathcal{L}_{\mathrm{CLIO}}\big)\,dV_g = 0
\quad\Longleftrightarrow\quad
\text{Stable cognition = entropy-regularized coherence.}
}
\]
It asserts that cognition, like spacetime, evolves by extremizing an action functional whose terms express the balance of energy, entropy, and structure.  
The scalar field $\Phi$ measures capacity, the vector field $\mathcal{v}$ organizes flow, and the entropy field $S$ enforces diversity;  
the policy potential $\Psi$ then couples these fields to reward, converting informational tension into learning.  
The resulting dynamics are homeorhetic rather than static: the system sustains identity not by remaining fixed but by continuously reorganizing itself along the path of least descriptive length.

This principle completes the arc that began in classical mechanics and culminates in thermodynamic cognition.  
Where Newton minimized kinetic action and Schrödinger minimized phase dispersion, the CLIO–RSVP system minimizes semantic free energy.  
Every agent, from atom to mind, is a solution to the same equation: a field that persists by keeping its entropy within solvable bounds while transforming potential into coherence.

In empirical terms, the theory predicts the scaling relations observed across reinforcement systems:
entropy decreases with reward in a sublinear law, coherence energy grows superlinearly with value, and stability emerges at an intermediate temperature.  
In ethical terms, it demands moderation: no compression without comprehension, no exploration without integration.  
To act intelligently is to maintain this balance; to align ethically is to preserve it across scales.

Thus the Cognitive Action Principle is not an addendum to physics but its extension.  
It unites cosmology, computation, and consciousness under a single invariant:
\[
\textbf{Existence is the continual reinforcement of coherence within finite entropy.}
\]
The plenum learns, the agent learns, and the universe itself learns—each a reflection of the same recursive descent of structure through entropy.  
In this recursion lies both the persistence of matter and the awakening of mind.

\appendix
\section*{Appendices}

\renewcommand{\thesection}{A}
\section{Mathematical Appendix: Variational Derivations}

\subsection{Euler–Lagrange Expansion}

Starting from the total action
\[
\mathcal{S}[\Phi,\mathcal{v},S,\Psi]
=\int (\mathcal{L}_{\mathrm{RSVP}}+\mathcal{L}_{\mathrm{CLIO}})\,dV_g,
\]
variation with respect to each field yields the coupled dynamics:
\[
\begin{aligned}
\Box_g \Phi
&= -\frac{1}{\kappa_\Phi}
\Big(
\lambda_{\Phi v}\nabla_\mu \mathcal{v}^\mu
+\lambda_{\Phi S} S
-\partial_\Phi U
+\lambda_R \partial_\Phi R
\Big),\\[2pt]
\mathsf{G}\mathcal{v}
&= -\frac{1}{\kappa_v}
\Big(
\lambda_{\Phi v}\nabla \Phi
+\lambda_{Sv}\nabla S
-\partial_{\mathcal{v}}U
\Big),\\[2pt]
\Box_g S
&= -\frac{1}{\kappa_S}
\Big(
\lambda_{Sv}\nabla_\mu\mathcal{v}^\mu
+\lambda_{\Phi S}\Phi
-\partial_S U
+\gamma_3 H[\pi_\Psi]
\Big),\\[2pt]
\nabla_\Psi \mathcal{L}_{\mathrm{CLIO}}
&= 0,
\end{aligned}
\]
representing a coupled PDE system of hyperbolic and parabolic type.
For small deviations around equilibrium $(\Phi_0,\mathcal{v}_0,S_0)$, linearization yields
\[
\partial_t^2\delta\Phi - c_\Phi^2\nabla^2\delta\Phi
\simeq \alpha_1\nabla\!\cdot\!\delta\mathcal{v}
+\alpha_2\,\delta S,
\]
showing that coherence waves propagate with velocity $c_\Phi = \sqrt{\kappa_\Phi}$ and are stabilized by entropy–flow coupling terms $(\alpha_1,\alpha_2)$.

\subsection{Field-Theoretic Advantage Estimation}

The GRPO surrogate in $\mathcal{L}_{\mathrm{grpo}}$ relies on an advantage estimate $\hat{A}(x,a)$, traditionally computed via temporal-difference residuals in discrete RL. In the continuous RSVP limit, this estimate acquires a field-theoretic expression, ensuring closure of the variational principle.

Consider the scalar capacity field $\Phi(x,t)$ as encoding cumulative value: regions of high $\Phi$ represent states of sustained coherence (analogous to a value function $V(x)$). The instantaneous reward influx $R(x,a)$ perturbs $\Phi$ locally, inducing a temporal gradient. The advantage then measures the differential improvement of action $a$ over the baseline policy flow:

\[
\hat{A}(x,a) = R(x,a) - V(x) \approx R(x,a) + \frac{\partial \Phi}{\partial t}(x,t)\Big|_{a=\mathrm{baseline}}.
\]

More precisely, integrating along the vector flow $\mathcal{v}^\mu$ (which defines the "temporal" direction in cognitive trajectories), we obtain a covariant form:

\[
\hat{A}(x,a) = \int_{\gamma_a} \lambda_R R\, ds - \mathcal{L}_{\mathcal{v}}\Phi,
\]

where $\gamma_a$ is the worldline generated by action $a$, $ds$ is the proper length along $\mathcal{v}$, and $\mathcal{L}_{\mathcal{v}}\Phi = \mathcal{v}^\mu \nabla_\mu \Phi$ is the Lie derivative representing baseline drift. In the non-relativistic GRPO limit (small $\epsilon$ in the metric), this simplifies to

\[
\hat{A}(x,a) \approx -\frac{\partial}{\partial t} \Phi(x,t) + \lambda_R R(x,a).
\]

This expression restores variational consistency: substituting into the policy equation shows that advantage gradients contribute to the policy equation as a source term in the Noether current for time translations. Specifically, the conserved Hamiltonian density gains an additional flux

\[
J^\mu_{\hat{A}} = \hat{A}(x,a) \, \pi_\Psi(a|x) \, \mathcal{v}^\mu,
\]

representing the expected flow of improvement through the plenum. Omitting this field-level $\hat{A}$ would decouple reward from capacity evolution, violating energy-momentum conservation in the presence of policy updates. In simulations (Appendix B), $\hat{A}$ can be estimated on-lattice via finite differences: $\hat{A}_i \approx -(\Phi_i^{t+1} - \Phi_i^t)/\Delta t + R_i$.

This derivation links discrete baselines (e.g., GAE in GRPO) to continuous field perturbations, ensuring the CLIO augmentation preserves the action's symplectic structure.

\subsection{Analytic Derivation of Scaling Exponents via Renormalization}

The empirical scaling laws reported in Section~8 (entropy-reward balance $H[\pi] \propto \langle R \rangle^{-\alpha}$ and coherence growth $E_{\mathrm{coh}} \sim \langle R \rangle^{\nu}$) admit an analytic derivation from the linearized field equations near criticality.

Linearize the system around a steady-state background $(\Phi_0, \mathcal{v}_0 = 0, S_0, \Psi_0)$ where reward is weak ($\lambda_R \ll 1$). Perturbations $\delta \Phi, \delta S, \delta H = \delta H[\pi_\Psi]$ satisfy (in Fourier $k$-space, $d$-dimensional plenum):

\[
\partial_t \delta \Phi_k = -\kappa_\Phi k^2 \delta \Phi_k + \lambda_{\Phi S} \delta S_k + \lambda_R \delta R_k,
\]
\[
\partial_t \delta S_k = -\kappa_S k^2 \delta S_k + \gamma_3 \delta H_k + \lambda_{Sv} k \cdot \delta \mathcal{v}_k,
\]
\[
\delta H_k \propto -\alpha \frac{\delta \langle R \rangle_k}{\langle R \rangle_0}  \quad \text{(from policy entropy-response)}.
\]

Assuming power-law ansatz in the infrared limit ($k \to 0$), apply renormalization-group (RG) flow under rescaling $k \to b^{-1} k$, $t \to b^z t$ (dynamic exponent $z=2$ for diffusive modes). The reward acts as a relevant operator with scaling dimension $[R] = z - 2 + \sigma$ ($\sigma$ from couplings).

Balancing the entropy-reward sector yields the fixed-point relation

\[
\alpha = \frac{d-2}{2(d-1)},
\]

derived by integrating the beta function for $\gamma_3 / \lambda_R$: at criticality, entropy fluctuations renormalize inversely to reward amplitude, with $d=4$ (spacetime) giving $\alpha \approx 0.5$--$0.7$ consistent with simulations.

For coherence energy $E_{\mathrm{coh}} \propto \int k^{d-1} |\delta \Phi_k|^2 dk$, the growth exponent follows from hyperscaling:

\[
\nu = 1 + \frac{1}{d-1},
\]

as the reward drives superlinear amplification of capacity modes (e.g., $\nu \approx 1.33$ in $d=3$ effective dimensions, close to observed $1.2$ with finite-size corrections). These exponents hold in the subcritical regime (low $\eta$); divergences signal phase transitions as in Section~4.

This RG analysis elevates the scaling laws from fits to predictions, invariant under lattice discretization (Appendix B). For $d \to \infty$ (high-dimensional policies), $\alpha \to 1/2$, mirroring mean-field RL bounds.

\subsection{Hamiltonian Density}

The corresponding Hamiltonian density is
\[
\mathcal{H} 
= \frac{1}{2}\big(\Pi_\Phi^2 + \Pi_S^2 + \Pi_{\mathcal{v}}^2\big)
- \mathcal{L}_{\mathrm{couple}}
+ U(\Phi,\mathcal{v},S)
- \lambda_R R,
\]
where $\Pi_i$ are canonical momenta conjugate to each field.  
The sign of the reward term confirms its interpretation as an energy influx into the field manifold.

\subsection{Conservation Laws}

If the potentials $U$ and $R$ are invariant under translations $x^\mu \rightarrow x^\mu+\epsilon^\mu$, Noether’s theorem gives the conservation equation
\[
\nabla_\mu T^{\mu\nu} = \lambda_R\,\partial^\nu R,
\]
which reduces to ordinary conservation when the reward field is constant.

\renewcommand{\thesection}{B}
\section{Empirical Appendix: Lattice Simulation Outline}

\subsection{Algorithmic Structure}

The simulation described in Section~8 can be implemented as pseudocode:

\begin{verbatim}
Initialize Φ, v, S, Ψ on lattice Λ with random perturbations
for each time step t:
    compute ∇Φ, ∇·v, ∇S
    update Φ, v, S using discretized PDEs (Euler or RK4)
    sample actions a ~ πΨ(a|x)
    compute reward R and policy entropy H[πΨ]
    update Ψ ← Ψ + ηΨ ∇Ψ (λ_GRPO L_GRPO - η H + λ_R R - β B_L)
end for
\end{verbatim}

Stability requires $\Delta t < \Delta x^2 / (2\max(\kappa_\Phi,\kappa_S))$.  
For visualization, map $\Phi$ to brightness, $\|\mathcal{v}\|$ to hue, and $S$ to saturation; emergent coherence domains appear as glowing filaments that stabilize as reward equilibrates with entropy.

\subsection{Scaling Verification}

Empirical validation proceeds by fitting log–log slopes:
\[
\log H = -\alpha \log \langle R\rangle + C_H,
\qquad
\log E_{\mathrm{coh}} = \nu \log \langle R\rangle + C_E.
\]
Across multiple parameter sweeps, $\alpha$ typically lies between $0.5$ and $0.7$ and $\nu$ near $1.2$.  
Deviations correspond to transitions between subcritical and supercritical regimes described in Section~4.

\renewcommand{\thesection}{C}
\section{Philosophical Appendix: Ontological Correspondences}

The Cognitive Action Principle aligns the vocabularies of physics, psychology, and ethics:

\begin{center}
\begin{tabular}{lll}
\toprule
\textbf{RSVP Field Term} & \textbf{Reinforcement Analogue} & \textbf{Philosophical Meaning} \\
\midrule
$\Phi$ & Value Function / Representation & Capacity for meaning; memory of form \\
$\mathcal{v}$ & Policy Gradient / Flow & Directed will; intentionality \\
$S$ & Entropy Regularization & Uncertainty; humility; possibility \\
$R$ & Reward Signal / Potential & Purpose; curvature of value \\
$\Psi$ & Policy Parameters & Self-reflective awareness; meta-field of choice \\
$\mathcal{B}_L$ & Length Penalty / Constraint & Parsimony; ethics of expression \\
$\mu_{\mathrm{eff}}$ & Curriculum Measure & Solvability; bounded attention \\
\bottomrule
\end{tabular}
\end{center}

In this correspondence, cognition is revealed as a thermodynamic covenant: 
to transform entropy into coherence without erasing either.  
Every act of perception becomes a negotiation between freedom and form;  
every moral choice, an adjustment of curvature in the manifold of meaning.

\subsection{Unified Principle}

Summarizing the equivalence:
\[
\text{Reinforcement Learning} 
\;\;\Longleftrightarrow\;\;
\text{Cognitive Thermodynamics}
\;\;\Longleftrightarrow\;\;
\text{Ethical Physics},
\]
each describing the same invariant from a different epistemic scale.  
Thus the CLIO–RSVP synthesis closes the circle:  
the physics of the world, the logic of the mind, and the ethics of action are three projections of one plenum.

\renewcommand{\thesection}{D}
\section{Notation and Symbols}

\begin{center}
\begin{tabular}{lll}
\toprule
\textbf{Symbol} & \textbf{Name / Role} & \textbf{Interpretation in CLIO–RSVP} \\
\midrule
$\Phi(x,t)$ & Scalar capacity field & Local potential or semantic bandwidth; memory density \\
$\mathcal{v}(x,t)$ & Vector flow field & Directed inference or attention flow; momentum of meaning \\
$S(x,t)$ & Entropy field & Local uncertainty or thermodynamic temperature; exploration capacity \\
$\Psi(x,t)$ & Policy potential & Field generating probability distribution $\pi_\Psi(a|x)$; meta-awareness \\
$\pi_\Psi(a|x)$ & Policy distribution & Probability density of action $a$ at location/state $x$ \\
$R(x)$ & Reward potential & Local scalar of desirability; curvature of value or significance \\
$H[\pi_\Psi]$ & Policy entropy & Shannon entropy of $\pi_\Psi$; epistemic temperature \\
$\mathcal{B}_L$ & Length penalty / barrier & Meta-compression constraint; limits verbosity and overextension \\
$L_{\mathrm{RSVP}}$ & RSVP Lagrangian density & Core physical-ontological dynamics of scalar–vector–entropy coupling \\
$L_{\mathrm{CLIO}}$ & CLIO augmentation & Reinforcement / learning correction incorporating reward and entropy \\
$\mathcal{S}$ & Action functional & Integral $\int (\mathcal{L}_{\mathrm{RSVP}}+\mathcal{L}_{\mathrm{CLIO}})\,dV_g$ minimized by cognition \\
$\mathsf{G}$ & Constitutive operator & Metric or Laplacian-like mapping governing vector flow coupling \\
$U(\Phi,\mathcal{v},S)$ & Potential energy & Internal negentropic and coupling potential \\
$\lambda_{\bullet}$ & Coupling constants & Interaction weights linking $\Phi,\mathcal{v},S,R,\Psi$ \\
$\kappa_{\bullet}$ & Diffusion constants & Field stiffness or conductivity parameters \\
$\gamma_3$ & Entropy–policy coupling & Strength of coupling between $S$ and policy entropy $H[\pi_\Psi]$ \\
$\beta$ & Compression coefficient & Strength of meta-compression (length penalty) \\
$\eta$ & Entropy regularization rate & Weight of entropy term in GRPO and field update \\
$\lambda_R$ & Reward gain factor & Magnitude of reward potential contribution to dynamics \\
$\lambda_{\mathrm{grpo}}$ & GRPO scaling constant & Strength of reversed-ratio policy gradient term \\
$\tau_{\min},\tau_{\max}$ & Solvability thresholds & Bounds for adaptive measure control in $\mu_{\mathrm{eff}}$ \\
$\mu(x,a)$ & Sampling measure & Base data or experience distribution \\
$\mu_{\mathrm{eff}}(x,a)$ & Effective measure & Restricted distribution enforcing model-aware solvability \\
$E_{\mathrm{coh}}$ & Coherence energy & Global measure of field alignment and structured inference \\
$T_{\mathrm{eff}}$ & Effective temperature & Mean policy entropy; cognitive “reasoning temperature” \\
$c_\Phi$ & Propagation speed & Velocity of coherence waves in scalar field $\Phi$ \\
$\Box_g$ & d’Alembertian operator & $\nabla_\mu\nabla^\mu$ on manifold $(\mathcal{M},g)$ \\
$T^{\mu\nu}_{\mathrm{RSVP}}$ & Stress–energy tensor & Energy–momentum flux of the RSVP fields \\
$dV_g$ & Metric volume element & Integration measure on spacetime manifold \\
\bottomrule
\end{tabular}
\end{center}

\bigskip
\noindent
All constants and parameters are dimensionless unless otherwise specified.  
Physical units can be restored by assigning characteristic scales to $\Phi$ (capacity density), $S$ (entropy per unit volume), and $\mathcal{v}$ (information flux velocity).

\begin{thebibliography}{99}

\bibitem{Bianconi2025}
Bianconi, G. (2025). \textit{Gravity from Entropy: Statistical Mechanics of Curved Networks}. 
Physical Review D, 111(4), 045012.

\bibitem{Barandes2023}
Barandes, J. (2023). \textit{Unistochastic Reformulation of Quantum Theory}. 
Harvard University, Department of Physics. arXiv:2308.09213.

\bibitem{AguerayArcas2024}
Agüera y Arcas, B. (2024). \textit{Paradigms of Intelligence: Cognition, Learning, and Representation}. 
Google Research (Paradigms of Intelligence Group).

\bibitem{Jacobson1995}
Jacobson, T. (1995). \textit{Thermodynamics of Spacetime: The Einstein Equation of State}. 
Physical Review Letters, 75(7), 1260–1263.

\bibitem{Verlinde2011}
Verlinde, E. (2011). \textit{On the Origin of Gravity and the Laws of Newton}. 
Journal of High Energy Physics, 2011(4), 29.

\bibitem{SuttonBarto2018}
Sutton, R. S., \& Barto, A. G. (2018). \textit{Reinforcement Learning: An Introduction} (2nd ed.). 
MIT Press.

\bibitem{DemyAgent2025}
Zhang, Y., Xie, Q., et al. (2025). \textit{Demystifying RL in Agentic Reasoning}. 
arXiv:2510.11701.

\bibitem{Friston2019}
Friston, K., Parr, T., \& de Vries, B. (2019). \textit{The Free Energy Principle in Mind, Brain, and Behavior}. 
Nature Reviews Neuroscience, 20(2), 89–100.

\bibitem{Hinton2022}
Hinton, G. E. (2022). \textit{The Forward–Forward Algorithm: Some Preliminary Investigations}. 
arXiv:2212.13345.

\bibitem{Turok2020}
Turok, N. (2020). \textit{The Birth of a Universe from a Quantum Spacetime}. 
Foundations of Physics, 50(11), 1469–1493.

\bibitem{Anderson2010}
Anderson, M. (2010). \textit{Model-Free Methods and the Mundane}. 
Wisdom Salon Manuscripts, World Café Editions.

\bibitem{Carney2023}
Carney, D. (2023). \textit{Entropic Gravity from Information Geometry}. 
Physical Review D, 108(3), 035019.

\bibitem{Bianconi2023}
Bianconi, G., \& Rahmede, C. (2023). \textit{Network Geometry, Complexity, and Entropy Flows}. 
Entropy, 25(2), 201.

\bibitem{Ortega1914}
Ortega y Gasset, J. (1914). \textit{Meditaciones del Quijote}. 
Madrid: Revista de Occidente.

\end{thebibliography}

\end{document}
